\documentclass[10pt,a4paper]{article}

\usepackage[T1]{fontenc}
\usepackage[utf8]{inputenc}
\usepackage{lmodern}
\usepackage{amsmath,amssymb,amsthm}
\usepackage{booktabs,array,tabularx}
\usepackage{geometry}
\usepackage{xcolor}
\usepackage{enumitem}
\usepackage{microtype}
\usepackage{hyperref}
\usepackage{fancyhdr}

\geometry{margin=1in}
\hypersetup{
  colorlinks=true,
  linkcolor=blue!55!black,
  citecolor=blue!55!black,
  urlcolor=blue!55!black,
  pdftitle={Riddle Me Zero: A One-Query Fresh-Message Forgery against Sphinx-in-the-Head},
  pdfauthor={Martin Feussner},
  pdfsubject={Cryptanalytic audit report for IACR ePrint 2024/649},
  pdfkeywords={group signatures, non-frameability, zero extension, Davies--Meyer, LowMC}
}
\setlist{nosep,leftmargin=*}
\setlength{\parindent}{1em}
\setlength{\parskip}{0.25em}
\setlength{\emergencystretch}{2em}
\Urlmuskip=0mu plus 1mu
\pagestyle{fancy}
\fancyhf{}
\fancyhead[L]{Riddle Me Zero}
\fancyhead[R]{24 October 2026}
\fancyfoot[C]{\thepage}
\renewcommand{\headrulewidth}{0.3pt}

\newtheorem{theorem}{Theorem}
\newcommand{\bits}{\{0,1\}}
\newcommand{\Hone}{H_1}
\newcommand{\sid}{\mathit{sid}}
\newcommand{\strr}{\mathit{str}}
\newcommand{\stt}{\mathit{stt}}
\newcommand{\com}{\mathit{com}}
\newcommand{\msg}{\mathit{msg}}
\newcommand{\GVf}{\mathsf{GVf}}
\newcommand{\GSig}{\mathsf{GSig}}
\newcommand{\Trace}{\mathsf{Trace}}
\newcommand{\TVf}{\mathsf{TVf}}

\title{\bfseries Riddle Me Zero:\\
A One-Query Fresh-Message Forgery against Sphinx-in-the-Head}
\author{Martin Feussner\\
\small Selmer Center, University of Bergen\\
\small \texttt{martin.feussner@uib.no}}
\date{24 October 2026}

\begin{document}
\maketitle

\noindent\fcolorbox{black!55}{black!4}{\parbox{0.94\textwidth}{
\textbf{AI provenance and review status.}\space
This audit used the OpenAI model Daybreak Blue at the ultra reasoning level.
OpenAI Codex generated and selected this paper's title and the
scheme-specific ``Riddle Me Zero'' wordplay. OpenAI Codex performed
AI-assisted target reconstruction, attack-history searching, cryptanalysis,
implementation coordination, experiment coordination, and technical writing.
A separate AI paper/profile review, hostile adversarial review, fresh-context
independent reproduction, and final exact-target attack-history recheck
passed. The exact public package was independently reviewed before release.
Independent human verification has not been performed. It is optional, and
there is no human-verification requirement.}}

\begin{abstract}
We analyze the current official 66-page version of
\emph{Sphinx-in-the-Head: Group Signatures from Symmetric Primitives}, IACR
ePrint 2024/649, frozen by SHA-256. The construction computes
$\sid=\Hone(\msg\mathbin\Vert\strr)$ and places $\sid$, rather than the two raw
operands, in its printed group-proof statement. At the frozen TEST2 profile,
Appendix~E.2 gives one- and two-block formulas satisfying
$\Hone(u)=\Hone(u\mathbin\Vert0^{255})$. After one honest signing query on the
empty bit string returns
$\Sigma_0=(u,\stt,\com,\pi_G)$, an adversary outputs the fresh 255-bit message
$u$ with $\Sigma^*=(0^{255},\stt,\com,\pi_G)$. Both pairs induce the same
$\sid$. Ordinary typed-paper $\GVf$ accepts, and a permitted tracing query
supplies a proof that $\TVf$ accepts for the honest member at epoch zero. This
wins the paper's Figure-8 non-frameability game and is also a one-query
adaptive fresh-message forgery.

A prospectively sealed, independently written standard-library model executed
the complete typed $\GSig/\GVf/\Trace/\TVf$ flow at
$(n,q,h,d,k,NR,NU)=(255,1024,1,16,70,1120,70)$. One independent trial
succeeded in 0.01501373597420752 seconds with 21,504~KiB peak
\texttt{RUSAGE\_SELF} and one worker; all eight controls rejected. This is a
deterministic algebraic transformation, so the single trial is not a
statistical success-rate estimate. The proof systems, credential service,
LowMC, M-FORS, and MPC-in-the-Head components are idealized. The official
HBGS repository is component-only and defines no outer message/randomizer
mapping, codec, full verifier, or relevant known-answer test. We therefore
make no author-byte, source-verifier, deployed-verifier, real-primitive, or
secret-recovery claim.
\end{abstract}

\section{Introduction}

Chen, Dong, Newton, and Wang propose Sphinx-in-the-Head (SITH), a fully
dynamic group signature built from symmetric primitives and an F-SPHINCS+
credential chain~\cite{chen2024}. A member signs an arbitrary finite bit string
using a fresh $n$-bit randomizer $\strr$. Both values are compressed to a
signature identifier $\sid=\Hone(\msg\Vert\strr)$ before the group proof is
formed. Verification, tracing, and tracing verification reconstruct that same
identifier.

Appendix~E.2 gives concrete LowMC-based cases for one, two, and three
$n$-bit blocks. In the one- and two-block cases, a zero second block is an
alias: feed-forward XOR with $0^n$ changes nothing. The signature randomizer
is public, so a signature on the empty message supplies exactly the block
needed to move those same $n$ bits from the randomizer field into the message
field. Replacing the randomizer by zero then preserves $\sid$ and the complete
printed proof statement.

This paper makes four contributions. First, it derives the zero-extension
identity directly from the target formulas. Second, it completes the formal
Figure-8 transcript, including the tracing query, exact-pair freshness,
$\tau=0$, and honest-member attribution. Third, it reports primary and
fresh-context independent full-dimension typed-paper executions, together
with public-only, mutation, revocation, and domain-separation controls.
Fourth, it states a strict source and implementation boundary: the result is
about the complete typed paper model, not an unprovided outer codec or a
deployed binary.

We claim no novelty for Davies--Meyer hashing, zero-extension observations,
fixed-point analysis, or message-randomizer retargeting in general. The
contribution is their concrete application to the printed
Sphinx-in-the-Head construction and its complete non-frameability endpoint.

\section{Target, profile, and claimed security}

The target is the official 66-page PDF of ePrint 2024/649, received on
28 April 2024. Its revision ordinal was not resolved, so the immutable target
is identified by the digest
\begin{center}
\small\ttfamily
67581ad5d1ae2455573cdb9b043344d91ca04788797318a7906ed500380a9881.
\end{center}
The paper presents post-quantum group signatures and describes the selected
row in a 128-bit quantum-security context. We record this as the authors'
claim, not as an independent contemporary security estimate.

\begin{table}[h]
\centering
\small
\begin{tabularx}{\textwidth}{@{}lX@{}}
\toprule
Item & Frozen value \\
\midrule
Profile & \texttt{SITH-2024-649-HBGS-TEST2-N255-Q1024-H1-}\newline
\texttt{D16-K70-NR1120-NU70-Q128-G2P10-PAPER-}\newline
\texttt{ALG2-ZEROPAD} \\
Core tuple & $(n,q,h,d,k)=(255,1024,1,16,70)$ \\
MPC schedule & $NR=1120$, $NU=70$, 16 parties \\
Group capacity & $2^{10}=1024$ members \\
M-FORS public value & $d k=1120$ bits \\
Top tree & 70 roots, 58 appended $0^{255}$ leaves, capacity 128, height 7 \\
Tree rule & Appendix-B Algorithm 2 zero padding; Appendix-E.5 lifting excluded \\
Security property & Dynamic group-signature non-frameability, Figure 8 \\
\bottomrule
\end{tabularx}
\caption{Frozen full-parameter typed paper profile.}
\label{tab:profile}
\end{table}

The paper states that its prototype implemented the credential chain and
group-proof signing/verification components, not the complete protocols.
The official HBGS repository is pinned at the following Git identity:
\begin{center}
\small\ttfamily
commit e13fe83cdc4a636dd8c72a37545090eef99812b7\\
tree 2cb34c40bff268dcdae3a1179ff6168f35fd91a4.
\end{center}
It is a component/test implementation~\cite{hbgs}. It exposes neither complete
outer $\mathsf{Join}$, $\GSig$, $\GVf$, $\Trace$, and $\mathsf{Revoke}$
interfaces nor a mapping from
its 64-byte \texttt{message\_digest} and 32-byte \texttt{nonce} buffers to
paper $\msg$, $\strr$, or $\sid$. It also supplies no complete outer codec,
arbitrary-length unary $\Hone$ dispatcher, SITH known-answer test, or full
verifier relevant to this transformation.

\section{The typed paper construction}

The group-signature syntax accepts $\msg\in\bits^*$. For TEST2 the signer
selects a 255-bit random string $\strr$ and computes
\begin{align}
\sid &= \Hone(\msg\mathbin\Vert\strr),\label{eq:sid}\\
\stt &= F(\mathit{tk}_u,\sid),\\
\mathit{sst} &= F(\mathit{sk}_u,\sid),\\
\mathit{gid} &= \Hone(\mathit{rpk}).
\end{align}
The remaining credential computation derives a commitment $\com$. The signer
proves the printed relation whose public tuple is
\begin{equation}
(\mathit{gp},\mathit{rpk},\mathit{gid},\sid,\stt,\com),
\label{eq:groupstatement}
\end{equation}
and emits
\begin{equation}
\Sigma=(\strr,\stt,\com,\pi_G).
\label{eq:signature}
\end{equation}
Raw $\msg$, raw $\strr$, and the revocation list are not separately listed in
the proof statement in Equation~\eqref{eq:groupstatement}.

Ordinary $\GVf$ parses Equation~\eqref{eq:signature}, recomputes
Equation~\eqref{eq:sid} and $\mathit{gid}$, rejects if some tracing key in the
supplied revocation list satisfies $\stt=F(\mathit{tk},\sid)$, and otherwise
verifies $\pi_G$ on Equation~\eqref{eq:groupstatement}. Ordinary $\Trace$
also recomputes $\sid$, finds the member whose tracing key matches $\stt$, and
issues $\pi_t$ for the public tracing tuple
\begin{equation}
(\mathit{id}_u,\mathit{gid},\stt,\sid).
\label{eq:tracestatement}
\end{equation}
The tracing verifier $\TVf$ evaluates that relation for its supplied message
and signature. Thus all four ordinary algorithms share the same typed $\sid$
boundary.

\section{The zero-extension identity}

The abstract construction selects
$\Hone:\bits^*\rightarrow\bits^n$. Appendix~E.2 gives the concrete one- and
two-block cases
\begin{align}
\mathsf{hash}_1(u)
  &=\mathsf{low\_mc}(u,0^n),\label{eq:oneblock}\\
\mathsf{hash}_{1a}(u,v)
  &=\mathsf{low\_mc}(u,v)\oplus v.\label{eq:twoblock}
\end{align}
When these cases instantiate the two relevant $\sid$ evaluations, substituting
$v=0^n$ gives
\begin{equation}
\Hone(u\mathbin\Vert0^n)
=\mathsf{low\_mc}(u,0^n)\oplus0^n
=\mathsf{low\_mc}(u,0^n)
=\Hone(u).
\label{eq:zeroextension}
\end{equation}

Equation~\eqref{eq:zeroextension} is not a cryptanalytic break of LowMC, a
collision search, or a block-cipher fixed-point computation. It is the literal
zero feed-forward case of the printed Davies--Meyer-style formulas. Nor do we
infer an unspecified framing rule for four or more blocks. The result needs
only the exact one- and two-block cases at $n=255$.

\section{One-query fresh-message transformation}

Let $\epsilon$ denote the unique zero-bit string. The adversary performs:
\begin{enumerate}
  \item Call $\mathsf{AddHU}$ to create the honest member.
  \item Query $\mathsf{GSigHU}(\epsilon)$ once, receiving
        $\Sigma_0=(u,\stt,\com,\pi_G)$ with $u\in\bits^{255}$.
  \item Set $\msg^*=u$ and
        $\Sigma^*=(0^{255},\stt,\com,\pi_G)$.
  \item Query $\mathsf{TraceU}(RL_0,\msg^*,\Sigma^*)$, obtaining the honest
        identity and tracing proof $\pi_t^*$.
  \item Output $(\msg^*,\Sigma^*,\pi_t^*,\tau=0)$.
\end{enumerate}
The signing and forged identifiers are
\begin{align}
\sid_0
 &=\Hone(\epsilon\mathbin\Vert u)
   =\Hone(u)
   =\mathsf{low\_mc}(u,0^{255}),\label{eq:sid0}\\
\sid_*
 &=\Hone(u\mathbin\Vert0^{255})
   =\mathsf{low\_mc}(u,0^{255})\oplus0^{255}
   =\sid_0.\label{eq:sidstar}
\end{align}
No secret is recovered. The transformation reads the public randomizer $u$
and moves it across the unframed concatenation boundary while retaining the
other three signature fields exactly.

\section{Correctness and formal consequence}

\begin{theorem}\label{thm:forgery}
Assume that Appendix~E.2's one- and two-block formulas instantiate the two
relevant $\sid$ evaluations in the frozen typed paper profile. From one honest
signature on $\epsilon$, the transformation above produces a signature
accepted on a fresh 255-bit message and wins Figure~8 at $\tau=0$. Ordinary
tracing and $\TVf$ attribute the output to the honest member.
\end{theorem}

\begin{proof}
Equations~\eqref{eq:sid0} and~\eqref{eq:sidstar} give $\sid_*=\sid_0$.
The root public key and $\mathit{gid}$ are global, while $\stt$, $\com$, and
$\pi_G$ are copied. Therefore every component of the printed public statement
in Equation~\eqref{eq:groupstatement} is unchanged. The empty list $RL_0$
contains no tracing key, so ordinary $\GVf(\msg^*,\Sigma^*,RL_0)$ reaches the
same proof verification and accepts.

The $\mathsf{TraceU}$ oracle first applies this ordinary verification and then
ordinary $\Trace$. Since the preserved $\stt$ still equals
$F(\mathit{tk}_h,\sid_*)$, tracing finds the honest identity. Its proof
$\pi_t^*$ authenticates the same tuple in Equation~\eqref{eq:tracestatement},
so the final $\TVf$ call accepts for that honest member.

The honest-response log contains only $(\epsilon,\Sigma_0)$. The forged
message has width 255, whereas $\epsilon$ has width zero. They are unequal even
if the numeric value sampled for $u$ is zero. Hence strict message freshness
holds, and the weaker Figure-8 requirement that the exact ordered pair
$(\msg^*,\Sigma^*)$ be absent from the log also holds. All predicates in
Figure~8 are therefore true.
\end{proof}

The tracing query is part of the non-frameability interface and supplies the
game-required $\pi_t^*$. It is not an additional signing query. A rule that
closed the game under reused tracing-token provenance rather than exact
$(\msg,\Sigma)$ pairs would define a different game from the printed one.

\section{Complexity and conditions}

The attack makes one honest signing query and one permitted tracing query. Its
own transformation copies three public fields and replaces the allocation of
255 public bits between $\msg$ and $\strr$; it requires no search, guessing,
secret information, credential fabrication, or proof fabrication. Algebraic
success is independent of the value of $u$. In particular, the probability
that $u=0^{255}$ creates no failure event because typed length still makes the
forged message fresh.

The condition is specification-level: the Appendix-E.2 one-/two-block cases
must govern the $\sid$ calls exactly as stated. A separately specified,
injective framing of the two operands would block this route, but no such
framing is present in the printed $\Hone(\msg\Vert\strr)$ expression or
proof statement.

\section{Full-profile reproduction}

Two independently developed source-independent typed models reached the same
endpoint. The corrected primary run used one trial and one worker, passed all
48 recorded checks, and took 3.48 seconds external wall time with 21,504~KiB
peak RSS. A fresh-context developer then wrote the public clean-room line
without inspecting or copying the delivered or primary attack code.

\begin{table}[h]
\centering
\small
\begin{tabular}{@{}lrrrl@{}}
\toprule
Run & Trials & Successes & Peak RSS & Time and scope \\
\midrule
Corrected primary model & 1 & 1 & 21,504 KiB & 3.48 s external wall \\
Independent clean-room v2 & 1 & 1 & 21,504 KiB & 0.01501373597420752 s internal \\
\bottomrule
\end{tabular}
\caption{Full-dimension typed-model executions. Neither time is a real-proof-system runtime.}
\label{tab:runs}
\end{table}

The exact independent implementation is 48,001 bytes with SHA-256
\begin{center}
\small\ttfamily
3b3b62d9018b28cda7a279fd2145697421f5b132a5d3728d6c99daf2178ee863,
\end{center}
and its retained 7,526-byte evidence JSON has SHA-256
\begin{center}
\small\ttfamily
f4b829edf9f822a8f44fbbf94108416a603571e7305d2e5bc1671e484995cfcd.
\end{center}
The evidence validator passed 166 of 166 assertions, and a separate hostile
result review found no blocker. The public program retains its historical
``prospectively sealed, not executed'' header so that its bytes remain exact;
the sealed run and code/evidence hash binding establish that those same bytes
were subsequently executed once.

The independent run reports one signing query on $\epsilon$, one tracing
query with $RL_0$, ordinary original and transformed acceptance, honest
tracing identity, final $\TVf$ acceptance, strict message freshness, and
Figure-8 return one. Both evaluated top trees used 70 real roots, 58 zero
leaves, height 7, and no Appendix-E.5 node lift. The observed $u$ was nonzero;
the zero edge in Theorem~\ref{thm:forgery} was established statically, not by
forcing that rare sample.

All eight precommitted controls passed:
\begin{enumerate}
  \item a public-only zero-query signature was rejected;
  \item public credential synthesis was rejected;
  \item public proof synthesis was rejected;
  \item proof reuse on a changed public statement was rejected;
  \item mutation of a signature field was rejected;
  \item replacement by a fixed nonzero $\strr$ was rejected;
  \item verification under a revocation list containing the honest tracing
        key was rejected; and
  \item a separate length-framed, domain-separated $\Hone$ world rejected the
        transformation while accepting its original signature.
\end{enumerate}

One success in one trial demonstrates the deterministic scoped transcript. It
is not an empirical estimate of a random success probability. The independent
internal interval ends after trial completion and before final JSON
serialization and \texttt{fsync}; it is not a whole-process benchmark.

\section{Evidence and implementation boundary}

The independent program evaluates the exact paper types, ordinary algorithm
flow, TEST2 dimensions, Figure-8 predicates, and Algorithm-2 top-tree shape.
To avoid vacuous public forgeries inside the model, credential issuance and
the two proof systems are stateful, hidden-key ideal capabilities. They verify
the printed witness relations before issuing statement-bound tags. Public
views can verify but cannot issue credentials or create proofs.

The implementation uses an independently domain-separated ideal surrogate for
LowMC, ideal M-FORS base values and credential-chain signatures, and a typed
fallback for $\Hone$ input shapes outside the printed one-, two-, and
three-block cases. It does not execute real LowMC, M-FORS, F-SPHINCS+, KKW,
Picnic, MPC-in-the-Head, Fiat--Shamir/NIZK, serialization, or parsing. The
attack equality itself uses only the exact one-/two-block formulas.

The official source commit and tree were stable, although GitHub archive
transport bytes differed between two acquisitions. That is retained as an
archive-byte qualification, not evidence of source-tree divergence. No
upstream source is included in the reproducer. Nothing here establishes
official-source interoperability, author-byte acceptance, a deployed
verifier, or a normative mapping for the source's external digest and nonce
buffers.

The evidence records are consistent with one authorized invocation, but the
absence of another artifact is not an operating-system audit of all historical
execution. These provenance limits do not affect the algebraic identity; they
bound the empirical claim.

\section{Related work and attack history}

The component formulas are Davies--Meyer-style block-cipher hashing as
standardized and discussed in generic hash literature~\cite{iso10118}. Prior
work studies fixed-message-length requirements for symmetric-key accumulator
and ring-signature constructions~\cite{derler2018}, randomized hashing in
digital signatures~\cite{gauravaram2009}, and generic fixed points of
block-cipher-based hashes~\cite{aumasson2009}. Those works are conceptual
background. They do not report this SITH transformation or its full
non-frameability endpoint, and we make no general technique-novelty claim.

The closest exact-paper criticism located was DGSP, ePrint 2025/760, which
discusses signature size and opening efficiency or scalability rather than a
forgery, recovery result, or non-frameability break. Generic results about
SPHINCS+, FORS, LowMC, Picnic, or MPC-in-the-Head concern different targets or
components.

No prior public substantive attack against the exact scheme, version, and
profile was located in the searches performed as of 2026-10-24. Confidence is
high for indexed exact-target literature and official/public repository
surfaces and medium-high globally. The search was bounded by HTTP 403 and WAF
blocks, GitHub rate limits, paywalls, disagreement and lag among citation
indexes, and the unavoidable possibility of private, deleted, embargoed, or
unindexed material. The statement is date-bounded rather than an absolute
priority claim.

\section{Repair direction}

An unambiguous, injective encoding should separate $\msg$ from $\strr$, for
example by domain separation and explicit bit lengths, and that same encoding
should be used consistently by $\GSig$, $\GVf$, $\Trace$, and $\TVf$. A
complete arbitrary-length $\Hone$ specification, outer codec, verifier, and
known-answer tests should cover the zero-length and block-boundary cases. The
raw typed operands, or a binding encoding of them, may also be incorporated
into the proof statement. This is a repair direction, not a proved drop-in
patch; a revised construction would require a new proof and interoperability
testing.

\section{Limitations}

\begin{enumerate}
  \item The result assumes the printed Appendix-E.2 one-/two-block formulas
        govern the two $\sid$ calls. It does not supply a missing general
        arbitrary-length dispatcher.
  \item The result is a source-independent typed-paper attack, not an
        official-source, byte-codec, interoperable, or deployed-verifier run.
  \item LowMC, M-FORS, F-SPHINCS+, MPC-in-the-Head, and both proof systems are
        not executed as real primitives; the proof and credential services
        are sound ideal capabilities.
  \item No secret key, equivalent key, credential, proof witness, or authority
        secret is recovered.
  \item One independent trial supports a deterministic transcript, not a
        statistical success-rate estimate. The $u=0^{255}$ edge is static,
        not empirically observed.
  \item The recorded internal runtime excludes final serialization and disk
        synchronization and is not an end-to-end author-implementation cost.
  \item Appendix-B Algorithm-2 zero padding is selected; no equivalence with
        Appendix-E.5 lifting is claimed.
  \item The official repository lacks the outer mapping and relevant KAT, and
        differing archive transport bytes prevent an archive-byte identity
        claim even though the Git commit and tree match.
  \item Artifact preservation supports, but cannot independently prove at the
        operating-system level, that no unrecorded invocation ever occurred.
  \item The attack-history conclusion is bounded by the recorded searches and
        their access and indexing limitations.
\end{enumerate}

\section{Reproducibility}

The public package retains the independent implementation and evidence
byte-for-byte. Quick mode validates existing evidence without importing the
model:
\begin{center}
\ttfamily ./reproducer/reproduce.sh --quick
\end{center}
A fresh post-publication user run requires an explicit nonexistent path:
\begin{center}
\ttfamily ./reproducer/reproduce.sh --full fresh-evidence.json
\end{center}
The latter starts a new random trial and then applies the standalone evidence
validator. The validator hashes the model as data; it does not import it. The
paper, scripts, evidence, reference outputs, and package checksums are included
with no external Python dependency.

\section{Conclusion}

The result is best described as a one-query adaptive fresh-message
non-frameability forgery in the complete typed paper model. It changes the
allocation of the same 255 bits between $\msg$ and $\strr$, leaves the proof
statement unchanged, and causes ordinary verification and tracing to accept
and blame the honest member. It is not an implementation-level LowMC,
M-FORS, MPC-in-the-Head, source-code, codec, or deployed-system claim.

\section*{Acknowledgements}
The computations were performed on the Norwegian Research and Education Cloud (NREC), using resources provided by the University of Bergen and the University of Oslo.

\begin{thebibliography}{9}

\bibitem{chen2024}
L.~Chen, C.~Dong, C.~J.~P. Newton, and Y.~Wang,
\newblock ``Sphinx-in-the-Head: Group Signatures from Symmetric Primitives,''
\newblock \emph{IACR Cryptology ePrint Archive}, Report 2024/649, 2024.
\newblock \url{https://eprint.iacr.org/2024/649}; doi:10.1145/3638763.

\bibitem{iso10118}
International Organization for Standardization,
\newblock \emph{ISO/IEC TR 10118-2:2010---Hash-functions using an $n$-bit
block cipher}, 2010.

\bibitem{derler2018}
D.~Derler, S.~Ramacher, and D.~Slamanig,
\newblock ``Post-Quantum Zero-Knowledge Proofs for Accumulators with
Applications to Ring Signatures from Symmetric-Key Primitives,''
\newblock in \emph{PQCrypto 2018}, LNCS 10786, pp.~419--440, 2018.
\newblock doi:10.1007/978-3-319-79063-3\_20; ePrint 2017/1154.

\bibitem{gauravaram2009}
P.~Gauravaram and L.~R. Knudsen,
\newblock ``On Randomizing Hash Functions to Strengthen the Security of
Digital Signatures,''
\newblock in \emph{EUROCRYPT 2009}, LNCS 5479, pp.~88--105, 2009.
\newblock doi:10.1007/978-3-642-01001-9\_5.

\bibitem{aumasson2009}
J.-P. Aumasson,
\newblock \emph{Analysis and Design of Symmetric Cryptographic Algorithms},
\newblock doctoral thesis, EPFL, 2009.

\bibitem{hbgs}
University of Surrey SCCS,
\newblock ``HBGS: test implementation for Sphinx-in-the-Head,''
\newblock commit \texttt{e13fe83cdc4a636dd8c72a37545090eef99812b7}.
\newblock \url{https://github.com/UoS-SCCS/HBGS}.

\end{thebibliography}

\end{document}
